\documentclass[aspectratio=169,10pt]{beamer}
\usetheme{Madrid}
\usecolortheme{dolphin}
\setbeamertemplate{navigation symbols}{}
\setbeamertemplate{caption}[numbered]

\usepackage[utf8]{inputenc}
\usepackage[T1]{fontenc}
\usepackage[spanish,es-noshorthands]{babel}
\usepackage{amsmath,amssymb,amsthm}
\usepackage{tikz}
\usepackage{pgfplots}
\pgfplotsset{compat=1.18}
\usetikzlibrary{shapes.geometric, arrows.meta, positioning, calc, trees, backgrounds}

\deftranslation[to=spanish]{Definition}{Definición}
\deftranslation[to=spanish]{Example}{Ejemplo}

\title[Prob. y Estadística]{Apéndice: Guía Visual de Gráficos Estadísticos}
\subtitle{40 Tipos de Gráficos para Estadística Descriptiva}
\institute[UNS]{Universidad Nacional del Sur \\ Departamento de Matemática}
\author{Probabilidad y Estadística (5765)}
\date{}

\begin{document}

\begin{frame}
\titlepage
\end{frame}

\begin{frame}{Contenidos del Apéndice}
\small
\begin{columns}[t]
\begin{column}{0.5\textwidth}
\begin{enumerate}
\item Gráfico de Barras Simple
\item Gráfico de Barras Agrupadas
\item Gráfico de Barras Apiladas
\item Gráfico de Barras 100\% Apiladas
\item Gráfico de Divergencia
\item Histogramas
\item Polígono de Frecuencias
\item Ojiva / Curva Acumulativa
\item Gráfico de Líneas
\item Gráfico de Áreas Apiladas
\item Gráfico de Dispersión
\item Gráfico de Burbujas
\item Gráfico de Sectores / Pie
\item Gráfico de Donas
\item Diagrama de Caja (Boxplot)
\item Gráfico de Violín
\item Diagrama de Tallo y Hojas
\item Mapa de Calor (Heatmap)
\item Gráfico de Radar / Araña
\item Gráfico Treemap
\end{enumerate}
\end{column}
\begin{column}{0.5\textwidth}
\begin{enumerate}
\setcounter{enumi}{20}
\item Gráfico de Sankey
\item Diagrama de Paradas / Waterfall
\item Gráfico de Pirámide Poblacional
\item Cartograma
\item Mapa Chopleth
\item Gráfico de Densidad (KDE)
\item Gráfico Q-Q Plot
\item Diagrama de Pareto
\item Gráfico Sparkline
\item Diagrama de Venn
\item Gráfico Correlograma / Matriz Correlación
\item Diagrama de Puntos (Dot Plot)
\item Diagrama de Voronoi
\item Diagrama de Cuerdas (Chord Diagram)
\item Diagrama de Red / Grafo
\item Gráfico de Pendientes (Slopechart)
\item Diagrama de Gantt
\item Nube de Palabras
\item Gráfico Candlestick (Velas Japonesas)
\item Diagrama de Mosaico (Mosaic Plot)
\end{enumerate}
\end{column}
\end{columns}
\end{frame}

% --- SLIDES 1 A 40 ---

\begin{frame}{1. Gráfico de Barras Simple}
\begin{columns}
\begin{column}{0.55\textwidth}
\begin{tikzpicture}
\begin{axis}[
    ybar, symbol=-1, bar width=15pt,
    symbolic x coords={A,B,C,D},
    xtick=data, ymin=0, ymax=100,
    width=7cm, height=5cm,
    ylabel={Porcentaje (\%)}
]
\addplot coordinates {(A,75) (B,45) (C,90) (D,30)};
\end{axis}
\end{tikzpicture}
\end{column}
\begin{column}{0.45\textwidth}
\begin{itemize}\small
\item \textbf{Tipo de variable:} Cualitativa nominal / ordinal o discreta.
\item \textbf{Caso de uso:} Comparar categorías discretas independientes.
\item \textbf{Contexto real:} Reportes de mercado sobre participación de sistemas operativos móviles (Android vs iOS).
\end{itemize}
\end{column}
\end{columns}
\end{frame}

\begin{frame}{2. Gráfico de Barras Agrupadas}
\begin{columns}
\begin{column}{0.55\textwidth}
\begin{tikzpicture}
\begin{axis}[
    ybar, bar width=8pt,
    symbolic x coords={2022,2023,2024},
    xtick=data, ymin=0, ymax=100,
    width=7cm, height=5cm, legend pos=north west
]
\addplot coordinates {(2022,40) (2023,50) (2024,65)};
\addplot coordinates {(2022,60) (2023,55) (2024,45)};
\legend{Prod A, Prod B}
\end{axis}
\end{tikzpicture}
\end{column}
\begin{column}{0.45\textwidth}
\begin{itemize}\small
\item \textbf{Tipo de variable:} Cualitativa bivariada o cualitativa + temporal.
\item \textbf{Caso de uso:} Comparar subgrupos a través de múltiples categorías.
\item \textbf{Contexto real:} Comparativa multianual de ventas por categoría de producto en informes financieros.
\end{itemize}
\end{column}
\end{columns}
\end{frame}

\begin{frame}{3. Gráfico de Barras Apiladas}
\begin{columns}
\begin{column}{0.55\textwidth}
\begin{tikzpicture}
\begin{axis}[
    ybar stacked, bar width=15pt,
    symbolic x coords={Q1,Q2,Q3,Q4},
    xtick=data, ymin=0, ymax=150,
    width=7cm, height=5cm
]
\addplot coordinates {(Q1,30) (Q2,45) (Q3,25) (Q4,60)};
\addplot coordinates {(Q1,40) (Q2,35) (Q3,50) (Q4,30)};
\end{axis}
\end{tikzpicture}
\end{column}
\begin{column}{0.45\textwidth}
\begin{itemize}\small
\item \textbf{Tipo de variable:} Cualitativa nominal + cuantitativa acumulativa.
\item \textbf{Caso de uso:} Mostrar composición interna del total por categorías.
\item \textbf{Contexto real:} Desglose del presupuesto nacional por ministerio y total consolidado.
\end{itemize}
\end{column}
\end{columns}
\end{frame}

\begin{frame}{4. Gráfico de Barras 100\% Apiladas}
\begin{columns}
\begin{column}{0.55\textwidth}
\begin{tikzpicture}
\begin{axis}[
    ybar stacked, bar width=15pt,
    symbolic x coords={Norte,Sur,Este},
    xtick=data, ymin=0, ymax=100,
    width=7cm, height=5cm
]
\addplot coordinates {(Norte,30) (Sur,60) (Este,45)};
\addplot coordinates {(Norte,70) (Sur,40) (Este,55)};
\end{axis}
\end{tikzpicture}
\end{column}
\begin{column}{0.45\textwidth}
\begin{itemize}\small
\item \textbf{Tipo de variable:} Cualitativa en proporciones.
\item \textbf{Caso de uso:} Comparar la proporción relativa de componentes sin importar el volumen total.
\item \textbf{Contexto real:} Porcentaje de energías renovables vs no renovables en la matriz de distintos países.
\end{itemize}
\end{column}
\end{columns}
\end{frame}

\begin{frame}{5. Gráfico de Divergencia}
\begin{columns}
\begin{column}{0.55\textwidth}
\begin{tikzpicture}
\begin{axis}[
    xbar, bar width=10pt,
    symbolic y coords={Item A, Item B, Item C},
    ytick=data, xmin=-50, xmax=50,
    width=7cm, height=5cm
]
\addplot coordinates {(-30,Item A) (-15,Item B) (-40,Item C)};
\addplot coordinates {(20,Item A) (35,Item B) (10,Item C)};
\end{axis}
\end{tikzpicture}
\end{column}
\begin{column}{0.45\textwidth}
\begin{itemize}\small
\item \textbf{Tipo de variable:} Cualitativa ordinal (escalas Likert).
\item \textbf{Caso de uso:} Visualizar acuerdo/desacuerdo o valores positivos/negativos respecto a un eje neutro.
\item \textbf{Contexto real:} Encuestas de satisfacción del cliente o intención de voto político.
\end{itemize}
\end{column}
\end{columns}
\end{frame}

\begin{frame}{6. Histograma}
\begin{columns}
\begin{column}{0.55\textwidth}
\begin{tikzpicture}
\begin{axis}[
    ybar interval,
    ymin=0, ymax=30,
    xmin=0, xmax=50,
    width=7cm, height=5cm,
    xlabel={Intervalo de Clase}, ylabel={Frecuencia}
]
\addplot coordinates {(0,5) (10,15) (20,28) (30,12) (40,4) (50,0)};
\end{axis}
\end{tikzpicture}
\end{column}
\begin{column}{0.45\textwidth}
\begin{itemize}\small
\item \textbf{Tipo de variable:} Cuantitativa continua.
\item \textbf{Caso de uso:} Analizar la distribución de frecuencia y forma (sesgo, simetría) de datos continuos.
\item \textbf{Contexto real:} Distribución de salarios o tiempos de procesamiento de solicitudes.
\end{itemize}
\end{column}
\end{columns}
\end{frame}

\begin{frame}{7. Polígono de Frecuencias}
\begin{columns}
\begin{column}{0.55\textwidth}
\begin{tikzpicture}
\begin{axis}[
    ymin=0, ymax=30,
    width=7cm, height=5cm,
    xlabel={Marcas de Clase}, ylabel={Frecuencia}
]
\addplot[thick, mark=*] coordinates {(5,5) (15,15) (25,28) (35,12) (45,4)};
\end{axis}
\end{tikzpicture}
\end{column}
\begin{column}{0.45\textwidth}
\begin{itemize}\small
\item \textbf{Tipo de variable:} Cuantitativa continua.
\item \textbf{Caso de uso:} Comparar la forma de dos o más distribuciones continuas en un mismo eje.
\item \textbf{Contexto real:} Distribuciones de notas en exámenes de distintos semestres.
\end{itemize}
\end{column}
\end{columns}
\end{frame}

\begin{frame}{8. Ojiva / Curva Acumulativa}
\begin{columns}
\begin{column}{0.55\textwidth}
\begin{tikzpicture}
\begin{axis}[
    ymin=0, ymax=100,
    width=7cm, height=5cm,
    xlabel={Límite Superior}, ylabel={Frecuencia Acumulada (\%)}
]
\addplot[thick, blue, mark=square*] coordinates {(0,0) (10,10) (20,35) (30,75) (40,92) (50,100)};
\end{axis}
\end{tikzpicture}
\end{column}
\begin{column}{0.45\textwidth}
\begin{itemize}\small
\item \textbf{Tipo de variable:} Cuantitativa continua acumulada.
\item \textbf{Caso de uso:} Determinar percentiles y proporción de datos debajo de un umbral.
\item \textbf{Contexto real:} Análisis de percentiles de crecimiento infantil (médico/pediátrico).
\end{itemize}
\end{column}
\end{columns}
\end{frame}

\begin{frame}{9. Gráfico de Líneas}
\begin{columns}
\begin{column}{0.55\textwidth}
\begin{tikzpicture}
\begin{axis}[
    width=7cm, height=5cm,
    xlabel={Tiempo}, ylabel={Valor}
]
\addplot[smooth, red, thick] coordinates {(1,10) (2,15) (3,12) (4,25) (5,30)};
\end{axis}
\end{tikzpicture}
\end{column}
\begin{column}{0.45\textwidth}
\begin{itemize}\small
\item \textbf{Tipo de variable:} Serie temporal / cuantitativa continua.
\item \textbf{Caso de uso:} Monitorear tendencias, ciclos o estacionalidad a lo largo del tiempo.
\item \textbf{Contexto real:} Evolución del IPC (índice de precios al consumidor) mensual publicado por el INDEC.
\end{itemize}
\end{column}
\end{columns}
\end{frame}

\begin{frame}{10. Gráfico de Áreas Apiladas}
\begin{columns}
\begin{column}{0.55\textwidth}
\begin{tikzpicture}
\begin{axis}[
    stack plots=y, area style,
    width=7cm, height=5cm,
    enlarge x limits=false
]
\addplot coordinates {(1,5) (2,10) (3,8) (4,12)} \closedcycle;
\addplot coordinates {(1,8) (2,12) (3,15) (4,18)} \closedcycle;
\end{axis}
\end{tikzpicture}
\end{column}
\begin{column}{0.45\textwidth}
\begin{itemize}\small
\item \textbf{Tipo de variable:} Serie temporal de componentes múltiples.
\item \textbf{Caso de uso:} Visualizar la contribución acumulada de varias partes al total a lo largo del tiempo.
\item \textbf{Contexto real:} Producción energética agregada por tipo de fuente (gas, hidroeléctrica, solar).
\end{itemize}
\end{column}
\end{columns}
\end{frame}

\begin{frame}{11. Gráfico de Dispersión (Scatter Plot)}
\begin{columns}
\begin{column}{0.55\textwidth}
\begin{tikzpicture}
\begin{axis}[
    only marks, mark=*, mark size=2pt,
    width=7cm, height=5cm,
    xlabel={Var X}, ylabel={Var Y}
]
\addplot coordinates {(1,2) (2,3.5) (3,3) (4,5) (5,5.5) (6,7)};
\end{axis}
\end{tikzpicture}
\end{column}
\begin{column}{0.45\textwidth}
\begin{itemize}\small
\item \textbf{Tipo de variable:} Dos variables cuantitativas continuas.
\item \textbf{Caso de uso:} Evaluar la relación, correlación o asociación entre dos variables cuantitativas.
\item \textbf{Contexto real:} Relación entre años de educación e ingresos promedio individuales.
\end{itemize}
\end{column}
\end{columns}
\end{frame}

\begin{frame}{12. Gráfico de Burbujas}
\begin{columns}
\begin{column}{0.55\textwidth}
\begin{tikzpicture}
\begin{axis}[
    width=7cm, height=5cm,
    xlabel={PIB per cápita}, ylabel={Esperanza de vida}
]
\node[circle, fill=blue!40, draw=blue, minimum size=15pt] at (axis cs:2,60) {};
\node[circle, fill=red!40, draw=red, minimum size=30pt] at (axis cs:5,75) {};
\node[circle, fill=green!40, draw=green, minimum size=10pt] at (axis cs:8,80) {};
\end{axis}
\end{tikzpicture}
\end{column}
\begin{column}{0.45\textwidth}
\begin{itemize}\small
\item \textbf{Tipo de variable:} Tres variables cuantitativas (Eje X, Eje Y, Tamaño).
\item \textbf{Caso de uso:} Representar datos trivariados simultáneamente en un plano bidimensional.
\item \textbf{Contexto real:} Visualizaciones socioeconómicas globales estilo Gapminder (Rosling).
\end{itemize}
\end{column}
\end{columns}
\end{frame}

\begin{frame}{13. Gráfico de Sectores / Pie}
\begin{columns}
\begin{column}{0.55\textwidth}
\begin{tikzpicture}
\draw[fill=blue!50] (0,0) -- (0:1.5) arc (0:120:1.5) -- cycle;
\draw[fill=red!50] (0,0) -- (120:1.5) arc (120:250:1.5) -- cycle;
\draw[fill=green!50] (0,0) -- (250:1.5) arc (250:360:1.5) -- cycle;
\end{tikzpicture}
\end{column}
\begin{column}{0.45\textwidth}
\begin{itemize}\small
\item \textbf{Tipo de variable:} Cualitativa nominal (pocas categorías, $\le 5$).
\item \textbf{Caso de uso:} Mostrar la proporción relativa de partes respecto al todo.
\item \textbf{Contexto real:} Distribución del mercado de navegadores web (Chrome, Safari, Firefox).
\end{itemize}
\end{column}
\end{columns}
\end{frame}

\begin{frame}{14. Gráfico de Donas}
\begin{columns}
\begin{column}{0.55\textwidth}
\begin{tikzpicture}
\draw[fill=blue!50] (0,0) -- (0:1.5) arc (0:140:1.5) -- cycle;
\draw[fill=orange!50] (0,0) -- (140:1.5) arc (140:360:1.5) -- cycle;
\fill[white] (0,0) circle (0.8);
\node at (0,0) {\small 100\%};
\end{tikzpicture}
\end{column}
\begin{column}{0.45\textwidth}
\begin{itemize}\small
\item \textbf{Tipo de variable:} Cualitativa nominal.
\item \textbf{Caso de uso:} Variación del gráfico de torta con espacio central reservado para un indicador clave (KPI).
\item \textbf{Contexto real:} Paneles de gestión empresarial (Dashboards) para metas de ventas.
\end{itemize}
\end{column}
\end{columns}
\end{frame}

\begin{frame}{15. Diagrama de Caja (Boxplot)}
\begin{columns}
\begin{column}{0.55\textwidth}
\begin{tikzpicture}
\begin{axis}[
    width=7cm, height=4cm,
    ytick={1}, yticklabels={Grupo A},
    xmin=0, xmax=100
]
\addplot+[boxplot prepared={lower whisker=10, lower quartile=30, median=50, upper quartile=70, upper whisker=90}] coordinates {};
\end{axis}
\end{tikzpicture}
\end{column}
\begin{column}{0.45\textwidth}
\begin{itemize}\small
\item \textbf{Tipo de variable:} Cuantitativa continua.
\item \textbf{Caso de uso:} Resumir la distribución mediante 5 números (Mín, Q1, Mediana, Q3, Máx) e identificar atípicos.
\item \textbf{Contexto real:} Comparación de rangos salariales por nivel educativo o industria.
\end{itemize}
\end{column}
\end{columns}
\end{frame}

\begin{frame}{16. Gráfico de Violín}
\begin{columns}
\begin{column}{0.55\textwidth}
\begin{tikzpicture}
\filldraw[fill=purple!30, draw=purple, thick] 
  (0,0) .. controls (0.5,1) and (0.8,2) .. (0,3)
  .. controls (-0.8,2) and (-0.5,1) .. (0,0);
\draw[black, thick] (0,0.5) -- (0,2.5);
\filldraw[black] (0,1.5) circle (2pt);
\end{tikzpicture}
\end{column}
\begin{column}{0.45\textwidth}
\begin{itemize}\small
\item \textbf{Tipo de variable:} Cuantitativa continua.
\item \textbf{Caso de uso:} Combina un Boxplot con un gráfico de densidad de probabilidad para ver multimodalidad.
\item \textbf{Contexto real:} Estudios genómicos y expresión de genes en bioinformática.
\end{itemize}
\end{column}
\end{columns}
\end{frame}

\begin{frame}{17. Diagrama de Tallo y Hojas}
\begin{columns}
\begin{column}{0.55\textwidth}
\begin{center}
\begin{tabular}{r|l}
\textbf{Tallo} & \textbf{Hojas} \\ \hline
1 & 2 3 5 8 \\
2 & 0 1 4 4 7 9 \\
3 & 2 5 8 \\
4 & 1 3
\end{tabular}
\end{center}
\end{column}
\begin{column}{0.45\textwidth}
\begin{itemize}\small
\item \textbf{Tipo de variable:} Cuantitativa discreta / continua.
\item \textbf{Caso de uso:} Mostrar la forma de la distribución reteniendo los valores numéricos exactos.
\item \textbf{Contexto real:} Análisis exploratorio de datos iniciales en muestras pequeñas.
\end{itemize}
\end{column}
\end{columns}
\end{frame}

\begin{frame}{18. Mapa de Calor (Heatmap)}
\begin{columns}
\begin{column}{0.55\textwidth}
\begin{tikzpicture}[scale=0.8]
\fill[blue!20] (0,0) rectangle (1,1);
\fill[blue!80] (1,0) rectangle (2,1);
\fill[blue!50] (0,1) rectangle (1,2);
\fill[blue!10] (1,1) rectangle (2,2);
\draw (0,0) grid (2,2);
\end{tikzpicture}
\end{column}
\begin{column}{0.45\textwidth}
\begin{itemize}\small
\item \textbf{Tipo de variable:} Bivariada continua / discreta codificada por intensidad de color.
\item \textbf{Caso de uso:} Visualizar patrones o concentraciones en matrices de datos densas.
\item \textbf{Contexto real:} Ocupación por horas e intensidades de tráfico en plataformas web.
\end{itemize}
\end{column}
\end{columns}
\end{frame}

\begin{frame}{19. Gráfico de Radar / Araña}
\begin{columns}
\begin{column}{0.55\textwidth}
\begin{tikzpicture}[scale=0.7]
\foreach \i in {1,...,5} {
  \draw[gray!50] (0,0) -- ({90-\i*72}:2);
}
\draw[gray!50] (18:1) -- (90:1) -- (162:1) -- (234:1) -- (306:1) -- cycle;
\draw[gray!50] (18:2) -- (90:2) -- (162:2) -- (234:2) -- (306:2) -- cycle;
\draw[red, thick, fill=red!20] (18:1.5) -- (90:1.8) -- (162:0.8) -- (234:1.2) -- (306:1.7) -- cycle;
\end{tikzpicture}
\end{column}
\begin{column}{0.45\textwidth}
\begin{itemize}\small
\item \textbf{Tipo de variable:} Multivariable cuantitativa continua (misma escala).
\item \textbf{Caso de uso:} Comparar perfiles multivariados de rendimiento o habilidades.
\item \textbf{Contexto real:} Evaluación de perfiles profesionales en recursos humanos o deportistas.
\end{itemize}
\end{column}
\end{columns}
\end{frame}

\begin{frame}{20. Gráfico Treemap}
\begin{columns}
\begin{column}{0.55\textwidth}
\begin{tikzpicture}[scale=0.8]
\draw[fill=blue!30] (0,0) rectangle (2.5,3);
\draw[fill=green!30] (2.5,0) rectangle (4,1.8);
\draw[fill=red!30] (2.5,1.8) rectangle (4,3);
\node at (1.25,1.5) {\small Cat A};
\node at (3.25,0.9) {\tiny Cat B};
\node at (3.25,2.4) {\tiny Cat C};
\end{tikzpicture}
\end{column}
\begin{column}{0.45\textwidth}
\begin{itemize}\small
\item \textbf{Tipo de variable:} Jerárquica / Cualitativa proporcional.
\item \textbf{Caso de uso:} Visualizar estructuras jerárquicas y tamaños relativos mediante rectángulos anidados.
\item \textbf{Contexto real:} Composición de exportaciones por sector económico nacional.
\end{itemize}
\end{column}
\end{columns}
\end{frame}

\begin{frame}{21. Gráfico de Sankey}
\begin{columns}
\begin{column}{0.55\textwidth}
\begin{tikzpicture}[scale=0.8]
\draw[fill=gray!30] (0,0) rectangle (0.5,3);
\draw[fill=blue!30] (3,1.5) rectangle (3.5,3);
\draw[fill=orange!30] (3,0) rectangle (3.5,1);
\fill[blue!20] (0.5,1.5) -- (3,1.5) -- (3,3) -- (0.5,3) -- cycle;
\fill[orange!20] (0.5,0) -- (3,0) -- (3,1) -- (0.5,1) -- cycle;
\end{tikzpicture}
\end{column}
\begin{column}{0.45\textwidth}
\begin{itemize}\small
\item \textbf{Tipo de variable:} Flujo entre estados continuos o categóricos.
\item \textbf{Caso de uso:} Representar transferencias, pérdidas o flujos entre etapas de un sistema.
\item \textbf{Contexto real:} Flujos de energía desde la generación primaria hasta el consumo final.
\end{itemize}
\end{column}
\end{columns}
\end{frame}

\begin{frame}{22. Diagrama de Cascadas / Waterfall}
\begin{columns}
\begin{column}{0.55\textwidth}
\begin{tikzpicture}
\begin{axis}[
    ybar, bar width=10pt,
    symbolic x coords={Inicio, Venta, Costo, Final},
    xtick=data, width=7cm, height=5cm
]
\addplot coordinates {(Inicio,100) (Venta,40) (Costo,-30) (Final,110)};
\end{axis}
\end{tikzpicture}
\end{column}
\begin{column}{0.45\textwidth}
\begin{itemize}\small
\item \textbf{Tipo de variable:} Cuantitativa acumulativa positiva/negativa.
\item \textbf{Caso de uso:} Explicar la transición entre un valor inicial y uno final por aportes y reducciones.
\item \textbf{Contexto real:} Estado de resultados financiero (EBITDA a Beneficio Neto).
\end{itemize}
\end{column}
\end{columns}
\end{frame}

\begin{frame}{23. Gráfico de Pirámide Poblacional}
\begin{columns}
\begin{column}{0.55\textwidth}
\begin{tikzpicture}
\begin{axis}[
    xbar stacked, bar width=8pt,
    symbolic y coords={60+, 30-59, 0-29},
    ytick=data, xmin=-50, xmax=50,
    width=7cm, height=5cm
]
\addplot coordinates {(-20,0-29) (-35,30-59) (-15,60+)};
\addplot coordinates {(22,0-29) (33,30-59) (18,60+)};
\end{axis}
\end{tikzpicture}
\end{column}
\begin{column}{0.45\textwidth}
\begin{itemize}\small
\item \textbf{Tipo de variable:} Demográfica (Edad por Sexo).
\item \textbf{Caso de uso:} Analizar la estructura por edad y sexo de una población.
\item \textbf{Contexto real:} Censos nacionales de población desarrollados por institutos de estadística.
\end{itemize}
\end{column}
\end{columns}
\end{frame}

\begin{frame}{24. Cartograma}
\begin{columns}
\begin{column}{0.55\textwidth}
\begin{tikzpicture}
\draw[fill=blue!30] (0,0) rectangle (2,2) node[midway]{\small Región A};
\draw[fill=green!30] (2,0) rectangle (2.8,2) node[midway]{\tiny B};
\draw[fill=red!30] (0,2) rectangle (2.8,3) node[midway]{\small Región C};
\end{tikzpicture}
\end{column}
\begin{column}{0.45\textwidth}
\begin{itemize}\small
\item \textbf{Tipo de variable:} Espacial geográfica modificada por variable cuantitativa.
\item \textbf{Caso de uso:} Distorsionar áreas geográficas según el valor de un indicador.
\item \textbf{Contexto real:} Mapas electorales distorsionados por número de electores por estado.
\end{itemize}
\end{column}
\end{columns}
\end{frame}

\begin{frame}{25. Mapa Coroplético (Choropleth Map)}
\begin{columns}
\begin{column}{0.55\textwidth}
\begin{tikzpicture}
\draw[fill=red!20] (0,0) -- (1,0) -- (0.8,1) -- (0,1) -- cycle;
\draw[fill=red!80] (1,0) -- (2,0) -- (2,1.2) -- (0.8,1) -- cycle;
\draw[fill=red!50] (0,1) -- (0.8,1) -- (1,2) -- (0,2) -- cycle;
\end{tikzpicture}
\end{column}
\begin{column}{0.45\textwidth}
\begin{itemize}\small
\item \textbf{Tipo de variable:} Espacial geográfica codificada por escala de color.
\item \textbf{Caso de uso:} Representar la variación de una tasa o densidad por regiones geográficas.
\item \textbf{Contexto real:} Mapas de tasa de incidencia de enfermedades o densidad poblacional.
\end{itemize}
\end{column}
\end{columns}
\end{frame}

\begin{frame}{26. Gráfico de Densidad (KDE)}
\begin{columns}
\begin{column}{0.55\textwidth}
\begin{tikzpicture}
\begin{axis}[
    width=7cm, height=5cm,
    xlabel={X}, ylabel={Densidad}
]
\addplot[smooth, blue, thick, fill=blue!10] coordinates {(-3,0) (-2,0.05) (-1,0.2) (0,0.4) (1,0.2) (2,0.05) (3,0)};
\end{axis}
\end{tikzpicture}
\end{column}
\begin{column}{0.45\textwidth}
\begin{itemize}\small
\item \textbf{Tipo de variable:} Cuantitativa continua estimando la función de densidad.
\item \textbf{Caso de uso:} Estimación suavizada y continua de la distribución subyacente de datos.
\item \textbf{Contexto real:} Análisis de modelos econométricos continuos.
\end{itemize}
\end{column}
\end{columns}
\end{frame}

\begin{frame}{27. Gráfico Q-Q (Quantile-Quantile Plot)}
\begin{columns}
\begin{column}{0.55\textwidth}
\begin{tikzpicture}
\begin{axis}[
    width=7cm, height=5cm,
    xlabel={Cuantiles Teóricos}, ylabel={Cuantiles Muestrales}
]
\addplot[red, dashed] coordinates {(-2,-2) (2,2)};
\addplot[only marks, mark=*] coordinates {(-1.8,-1.7) (-1,-0.9) (0,0.1) (1,0.95) (1.8,2.1)};
\end{axis}
\end{tikzpicture}
\end{column}
\begin{column}{0.45\textwidth}
\begin{itemize}\small
\item \textbf{Tipo de variable:} Cuantitativa (evaluación diagnóstica).
\item \textbf{Caso de uso:} Comparar la distribución muestral con una distribución teórica (ej. Normal).
\item \textbf{Contexto real:} Verificación de supuestos de normalidad en residuos de regresión.
\end{itemize}
\end{column}
\end{columns}
\end{frame}

\begin{frame}{28. Diagrama de Pareto}
\begin{columns}
\begin{column}{0.55\textwidth}
\begin{tikzpicture}
\begin{axis}[
    ybar, bar width=10pt,
    width=7cm, height=5cm,
    symbolic x coords={A,B,C,D}, xtick=data
]
\addplot coordinates {(A,50) (B,30) (C,12) (D,8)};
\addplot[sharp plot, red, mark=*] coordinates {(A,50) (B,80) (C,92) (D,100)};
\end{axis}
\end{tikzpicture}
\end{column}
\begin{column}{0.45\textwidth}
\begin{itemize}\small
\item \textbf{Tipo de variable:} Cualitativa nominal + cuantitativa acumulada.
\item \textbf{Caso de uso:} Identificar las causas principales responsables de la mayoría de los efectos (regla 80/20).
\item \textbf{Contexto real:} Control de calidad industrial y gestión de defectos de fabricación.
\end{itemize}
\end{column}
\end{columns}
\end{frame}

\begin{frame}{29. Gráfico Sparkline}
\begin{columns}
\begin{column}{0.55\textwidth}
\vspace{1cm}
Texto en reporte: Tendencia reciente \begin{tikzpicture}\draw[thick, green!60!black] (0,0) -- (0.3,0.2) -- (0.6,-0.1) -- (0.9,0.4) -- (1.2,0.3);\end{tikzpicture} positiva.
\end{column}
\begin{column}{0.45\textwidth}
\begin{itemize}\small
\item \textbf{Tipo de variable:} Serie temporal condensada.
\item \textbf{Caso de uso:} Mostrar variaciones rápidas e inline dentro de un texto o tabla de datos.
\item \textbf{Contexto real:} Tablas resumen de mercados bursátiles en periódicos financieros.
\end{itemize}
\end{column}
\end{columns}
\end{frame}

\begin{frame}{30. Diagrama de Venn}
\begin{columns}
\begin{column}{0.55\textwidth}
\begin{tikzpicture}[scale=0.8]
\draw[fill=blue!30, opacity=0.5] (-0.5,0) circle (1.2);
\draw[fill=red!30, opacity=0.5] (0.5,0) circle (1.2);
\node at (-1.2,0) {A};
\node at (1.2,0) {B};
\node at (0,0) {\tiny $A\cap B$};
\end{tikzpicture}
\end{column}
\begin{column}{0.45\textwidth}
\begin{itemize}\small
\item \textbf{Tipo de variable:} Conjuntos de datos categóricos.
\item \textbf{Caso de uso:} Mostrar las relaciones lógicas de solapamiento e intersección entre grupos.
\item \textbf{Contexto real:} Análisis de solapamiento de audiencias entre plataformas digitales.
\end{itemize}
\end{column}
\end{columns}
\end{frame}

\begin{frame}{31. Gráfico Correlograma}
\begin{columns}
\begin{column}{0.55\textwidth}
\begin{tikzpicture}[scale=0.8]
\draw[fill=blue!90] (0,1) circle (0.3);
\draw[fill=blue!20] (1,1) circle (0.1);
\draw[fill=red!70] (0,0) circle (0.25);
\draw[fill=blue!90] (1,0) circle (0.3);
\draw (0,-0.5) grid (1,1.5);
\end{tikzpicture}
\end{column}
\begin{column}{0.45\textwidth}
\begin{itemize}\small
\item \textbf{Tipo de variable:} Matriz multivariada de correlaciones.
\item \textbf{Caso de uso:} Visualizar la magnitud y dirección de correlaciones entre múltiples variables.
\item \textbf{Contexto real:} Análisis multivariado preliminar en ciencia de datos.
\end{itemize}
\end{column}
\end{columns}
\end{frame}

\begin{frame}{32. Diagrama de Puntos (Dot Plot)}
\begin{columns}
\begin{column}{0.55\textwidth}
\begin{tikzpicture}
\begin{axis}[
    only marks, mark=*,
    symbolic y coords={G1,G2,G3},
    ytick=data, xmin=0, xmax=100,
    width=7cm, height=4cm
]
\addplot coordinates {(20,G1) (50,G1) (80,G1)};
\addplot coordinates {(30,G2) (60,G2)};
\addplot coordinates {(40,G3) (70,G3) (90,G3)};
\end{axis}
\end{tikzpicture}
\end{column}
\begin{column}{0.45\textwidth}
\begin{itemize}\small
\item \textbf{Tipo de variable:} Cuantitativa discreta / continua agrupada.
\item \textbf{Caso de uso:} Mostrar observaciones individuales en un eje unidimensional para muestras pequeñas.
\item \textbf{Contexto real:} Comparativa de puntajes de pruebas estandarizadas entre instituciones.
\end{itemize}
\end{column}
\end{columns}
\end{frame}

\begin{frame}{33. Diagrama de Voronoi}
\begin{columns}
\begin{column}{0.55\textwidth}
\begin{tikzpicture}[scale=0.7]
\draw (0,0) -- (1,2) -- (2.5,1.5) -- (2,0) -- cycle;
\draw (2.5,1.5) -- (3.5,2.5) -- (4,1) -- (2,0);
\filldraw (0.8,0.8) circle (2pt);
\filldraw (2.2,1) circle (2pt);
\filldraw (3.2,1.8) circle (2pt);
\end{tikzpicture}
\end{column}
\begin{column}{0.45\textwidth}
\begin{itemize}\small
\item \textbf{Tipo de variable:} Espacial continua (partición del plano).
\item \textbf{Caso de uso:} Determinar regiones de influencia basadas en la distancia al punto más cercano.
\item \textbf{Contexto real:} Delimitación de zonas de cobertura de hospitales o antenas celulares.
\end{itemize}
\end{column}
\end{columns}
\end{frame}

\begin{frame}{34. Diagrama de Cuerdas (Chord Diagram)}
\begin{columns}
\begin{column}{0.55\textwidth}
\begin{tikzpicture}[scale=0.8]
\draw (0,0) circle (1.5);
\draw[blue, thick] (50:1.5) to[out=230,in=310] (130:1.5);
\draw[red, thick] (200:1.5) to[out=20,in=110] (320:1.5);
\end{tikzpicture}
\end{column}
\begin{column}{0.45\textwidth}
\begin{itemize}\small
\item \textbf{Tipo de variable:} Interrelaciones entre entidades categóricas.
\item \textbf{Caso de uso:} Visualizar flujos bidireccionales o relaciones de origen-destino.
\item \textbf{Contexto real:} Flujos de migración global entre continentes.
\end{itemize}
\end{column}
\end{columns}
\end{frame}

\begin{frame}{35. Diagrama de Red / Grafo}
\begin{columns}
\begin{column}{0.55\textwidth}
\begin{tikzpicture}[node distance=1.2cm, main/.style = {draw, circle, fill=blue!20}]
\node[main] (1) {A};
\node[main] (2) [above right of=1] {B};
\node[main] (3) [below right of=1] {C};
\draw (1) -- (2); \draw (1) -- (3); \draw (2) -- (3);
\end{tikzpicture}
\end{column}
\begin{column}{0.45\textwidth}
\begin{itemize}\small
\item \textbf{Tipo de variable:} Relacional / Estructura de red.
\item \textbf{Caso de uso:} Representar relaciones, conexiones y nodos en un sistema complejo.
\item \textbf{Contexto real:} Análisis de redes sociales e interconexiones de usuarios.
\end{itemize}
\end{column}
\end{columns}
\end{frame}

\begin{frame}{36. Gráfico de Pendientes (Slopechart)}
\begin{columns}
\begin{column}{0.55\textwidth}
\begin{tikzpicture}
\draw[gray] (0,0) -- (0,3) node[above]{2020};
\draw[gray] (2,0) -- (2,3) node[above]{2024};
\draw[red, thick] (0,2.5) -- (2,1.0) node[right]{\tiny País A};
\draw[blue, thick] (0,0.8) -- (2,2.2) node[right]{\tiny País B};
\end{tikzpicture}
\end{column}
\begin{column}{0.45\textwidth}
\begin{itemize}\small
\item \textbf{Tipo de variable:} Cuantitativa comparativa en dos puntos en el tiempo.
\item \textbf{Caso de uso:} Evaluar cambios en rangos o valores entre dos momentos específicos.
\item \textbf{Contexto real:} Cambios en la clasificación del ranking del IDH de países.
\end{itemize}
\end{column}
\end{columns}
\end{frame}

\begin{frame}{37. Diagrama de Gantt}
\begin{columns}
\begin{column}{0.55\textwidth}
\begin{tikzpicture}
\draw (0,1.5) node[left]{\tiny Tarea 1} -- (3,1.5);
\fill[blue!60] (0,1.3) rectangle (1.5,1.7);
\draw (0,0.8) node[left]{\tiny Tarea 2} -- (3,0.8);
\fill[orange!60] (1.2,0.6) rectangle (2.8,1.0);
\end{tikzpicture}
\end{column}
\begin{column}{0.45\textwidth}
\begin{itemize}\small
\item \textbf{Tipo de variable:} Temporal / Planificación de proyectos.
\item \textbf{Caso de uso:} Ilustrar el cronograma y la dependencia temporal de actividades.
\item \textbf{Contexto real:} Gestión y seguimiento de proyectos de infraestructura.
\end{itemize}
\end{column}
\end{columns}
\end{frame}

\begin{frame}{38. Nube de Palabras (Word Cloud)}
\begin{columns}
\begin{column}{0.55\textwidth}
\begin{center}
{\Huge \textbf{Datos}} {\large Estadística} \\[0.2cm]
{\small Muestra} {\Huge \textbf{Probabilidad}}
\end{center}
\end{column}
\begin{column}{0.45\textwidth}
\begin{itemize}\small
\item \textbf{Tipo de variable:} Frecuencia de texto cualitativo.
\item \textbf{Caso de uso:} Representar de forma visual la frecuencia relativa de palabras en un cuerpo de texto.
\item \textbf{Contexto real:} Análisis de menciones en redes sociales y discursos políticos.
\end{itemize}
\end{column}
\end{columns}
\end{frame}

\begin{frame}{39. Gráfico Candlestick (Velas Japonesas)}
\begin{columns}
\begin{column}{0.55\textwidth}
\begin{tikzpicture}
\draw (0.5,0.5) -- (0.5,2.5);
\draw[fill=green!60] (0.3,1.0) rectangle (0.7,2.0);
\draw (1.5,0.8) -- (1.5,2.8);
\draw[fill=red!60] (1.3,1.2) rectangle (1.7,2.4);
\end{tikzpicture}
\end{column}
\begin{column}{0.45\textwidth}
\begin{itemize}\small
\item \textbf{Tipo de variable:} Cuantitativa de series financieras (Apertura, Máximo, Mínimo, Cierre).
\item \textbf{Caso de uso:} Describir el movimiento de precios de activos financieros en un intervalo.
\item \textbf{Contexto real:} Plataformas de trading y análisis de acciones bursátiles.
\end{itemize}
\end{column}
\end{columns}
\end{frame}

\begin{frame}{40. Diagrama de Mosaico (Mosaic Plot)}
\begin{columns}
\begin{column}{0.55\textwidth}
\begin{tikzpicture}[scale=0.8]
\draw[fill=blue!40] (0,0) rectangle (1.5,2);
\draw[fill=blue!20] (0,2) rectangle (1.5,3);
\draw[fill=red!40] (1.5,0) rectangle (3,1);
\draw[fill=red!20] (1.5,1) rectangle (3,3);
\end{tikzpicture}
\end{column}
\begin{column}{0.45\textwidth}
\begin{itemize}\small
\item \textbf{Tipo de variable:} Tablas de contingencia multivariadas.
\item \textbf{Caso de uso:} Examinar relaciones entre dos o más variables categóricas simultáneamente.
\item \textbf{Contexto real:} Estudios epidemiológicos multivariable sobre factores de riesgo.
\end{itemize}
\end{column}
\end{columns}
\end{frame}

\end{document}
